\documentclass[tikz,border=2mm]{standalone}
\usetikzlibrary{math}

\makeatletter
\newif\iffounddigit

% 1. 递归检查每一个字符
\def\scan@for@digits#1{%
  \ifx#1\@empty
    % 扫描结束
  \else
    % 针对 0-9 每一个数字进行硬比对
    \if#10\global\founddigittrue\fi
    \if#11\global\founddigittrue\fi
    \if#12\global\founddigittrue\fi
    \if#13\global\founddigittrue\fi
    \if#14\global\founddigittrue\fi
    \if#15\global\founddigittrue\fi
    \if#16\global\founddigittrue\fi
    \if#17\global\founddigittrue\fi
    \if#18\global\founddigittrue\fi
    \if#19\global\founddigittrue\fi
    % 继续处理下一个 token
    \expandafter\scan@for@digits
  \fi
}

% 2. 预处理：将内容转为字符串并剔除“\”的影响
\newcommand{\checkfornumbers}[1]{%
  \global\founddigitfalse
  % 将参数转为纯文本字符串
  \edef\temp@str{\detokenize{#1}}%
  % 启动递归扫描
  \expandafter\scan@for@digits\temp@str\@empty
}

\tikzset{
  define/@circle-point/.code args={#1,#2}{%
  \message{^^J[PASS] \detokenize{#2} is a Point^^J}%
},
define/@circle-radius/.code args={#1,#2}{%
  \message{^^J[MATCH] \detokenize{#2} contains Numbers/Expression^^J}%
},
  define/circle/.code args={#1,#2}{%
    \checkfornumbers{#2}%
    \iffounddigit
      \pgfkeys{/tikz/define/@circle-radius={#1,#2}}%
    \else
      \pgfkeys{/tikz/define/@circle-point={#1,#2}}%
    \fi
  },
}
\makeatother

\begin{document}
\begin{tikzpicture}
  \tikzmath{
    \O{1} = 1; \O{2} = 1;
    \P{1} = 2; \P{2} = 3;
    \radius = 5;
  }

  % 情况 A：只有 \P，不含数字 0-9，应该识别为 Point
  \tikzset{define/circle={\O,\P}}

  % 情况 B：含数字 1，应该识别为 Radius
  \tikzset{define/circle={\O,\P{1}}}

  % 情况 C：含常数，应该识别为 Radius
  \tikzset{define/circle={\O, 5cm}}

  \tikzset{define/circle={\O, \radius}}%failed

\end{tikzpicture}
\end{document}

% founddigitfalse 要设置为全局吗
% **必须设置为全局（Global）。**

% 在 TikZ 的 `pgfkeys` 环境中，`.code` 或 `.code args` 执行的代码块实际上是在一个**隐藏的分组（Group/Scope）**中运行的。

% ### 为什么要用全局？

% 如果你在 `.code` 块中使用 `\founddigitfalse`（局部赋值）：
% 1.  进入 `define/circle/.code`，开启一个 LaTeX 分组。
% 2.  执行 `\checkfornumbers`，变量被设为 `true` 或 `false`。
% 3.  **退出代码块时，分组结束，变量的值会立即回滚（回滚到执行该键之前的状态）。**

% 这会导致你在外面（或者接下来的逻辑判断中）永远拿不到正确的更新结果。

% ### 改进后的标准化写法

% 为了确保安全，你应该使用 `\global\founddigitfalse` 或 `\founddigitfalse` 的全局版本（如果使用的是原生 `\newif`，通常配合 `\global` 使用）。

% 在 LaTeX 中，`\newif\iffound` 会自动生成三个宏：
% *   `\foundtrue`
% *   `\foundfalse`
% *   `\iffound`

% 当你需要全局操作时，标准的做法是：

% ```latex
% \newcommand{\checkfornumbers}[1]{%
%   \global\founddigitfalse % 初始化为全局假
%   \edef\temp@str{\detokenize{#1}}%
%   \expandafter\scan@for@digits\temp@str\@empty
% }

% \def\scan@for@digits#1{%
%   \ifx#1\@empty\else
%     % 如果匹配到数字，必须全局设为真
%     \if#10\global\founddigittrue\fi
%     \if#11\global\founddigittrue\fi
%     \if#12\global\founddigittrue\fi
%     \if#13\global\founddigittrue\fi
%     \if#14\global\founddigittrue\fi
%     \if#15\global\founddigittrue\fi
%     \if#16\global\founddigittrue\fi
%     \if#17\global\founddigittrue\fi
%     \if#18\global\founddigittrue\fi
%     \if#19\global\founddigittrue\fi
%     \expandafter\scan@for@digits
%   \fi
% }
% ```

% ### 为什么你之前“全都识别成一种结果”可能与此有关？
% 如果 `\founddigit` 默认是 `false`，且你没有使用 `\global\founddigittrue`，那么即便程序在扫描时发现了数字，一旦退出 `\scan@for@digits` 所在的那个小环境，它就又变回 `false` 了。结果就是：**永远进不了 Number 分支**。

% ### 总结
% 1.  **初始化**：`\global\founddigitfalse`
% 2.  **命中目标**：`\global\founddigittrue`
% 3.  **判断**：`\iffounddigit`（判断不需要加 global）

% 这样就能确保你在 `\tikzset` 的复杂逻辑嵌套里，依然能准确捕获到判断结果。